% After training the transformer, we extract a \strips model $M'$ for use with an off-the-shelf symbolic planner, removing the auxiliary \texttt{init-start}, $\texttt{init-}p$, and $\texttt{goal-}p$ actions. The extraction procedure depends on the supervision setting.

After training, we extract a lifted \strips model $M'$ and discard the auxiliary \texttt{init-start}, $\texttt{init-}p$, and $\texttt{goal-}p$ schemas.
The extraction procedure depends on the supervision setting.

\paragraph{State reconstruction.}
Under final-state supervision, both Transformers are trained to answer ground-atom queries about the state reached by a trace. After training, we apply these queries to action-sequence prefixes to infer unobserved intermediate-state values. For each prefix ending in a domain action $a_i$, we query all candidate atoms whose arguments are drawn from the objects in $a_i$ and, when present, $a_{i+1}$, together with all nullary atoms. These are the only atoms that can be effects of $a_i$ or preconditions of $a_{i+1}$ under the \strips representation considered here. The observed initial state and these predictions provide the local pre- and
post-state observations needed for every action occurrence.

For an occurrence $a_i=\alpha(\mathbf{o}_i)$, a queried atom
$p(\mathbf{o}_i^\rho)$ is lifted to the candidate atom represented by
$(p,\rho)$.
To tolerate prediction errors, we include a candidate as a precondition when it
is predicted true before at least $95\%$ of the occurrences of $\alpha$.
For effects, every occurrence votes for \emph{add}, \emph{delete}, or
\emph{none} according to its predicted transition, and we assign the modal
label.
A modal label of \emph{none} yields no effect.
The resulting preconditions and effects define $M'$.
Although the LS Transformer also admits direct parameter extraction, we use
this common architecture-independent procedure for both models under
final-state supervision.

%\paragraph{State reconstruction.}
% Under final-state supervision, both the LS and LB Transformers are trained to predict the final state $s_F$ of each trace. We exploit this capability after training to reconstruct the unobserved intermediate states of the training traces, in a procedure that we call \emph{state reconstruction}. In particular, for every prefix ending in a domain action $a_i$, we query the trained transformer using $\texttt{goal-}p$ actions to determine the truth values of ground atoms at the state reached after $a_i$. 
%For every pair of consecutive trace actions $a_i$ and $a_{i+1}$, we only query atoms instantiated over objects appearing in either action, since these are the only atoms that can occur as effects of $a_i$ or preconditions of $a_{i+1}$. 
%This yields reconstructed state-action traces $s_0,a_0,s_1,a_1,\ldots,a_{n-1},s_n.$

% We then lift the ground atoms relative to the arguments of the corresponding action schema and aggregate the resulting observations across the training traces. A lifted atom is included as a precondition when it is predicted true before at least $95\%$ of the occurrences of the action schema. To infer effects, we compare the reconstructed states before and after each action occurrence: an atom votes for \emph{add} if it changes from false to true, \emph{delete} if it changes from true to false, and \emph{none} otherwise. The most frequent label determines its effect.
%The resulting preconditions and effects define the extracted lifted \strips model. This procedure is independent of the transformer architecture and is therefore used for both the LS and LB Transformers under final-state supervision.

\paragraph{Direct extraction via symbolic alignment.}
Under applicability supervision, no successor-state queries are available, so state reconstruction cannot be used. For the LS Transformer, we instead extract $M'$ directly from its binarized parameters:
$q_{p,\alpha,\rho}$, $k_{p,\alpha,\rho}$, and
$v_{p,\alpha,\rho}$ specify precondition membership, effect membership, and effect sign, respectively. The fixed $\texttt{init-}p$ schemas align hidden predicate heads with the corresponding predicates $p$. Heads without such an alignment retain anonymous names. Because the LB parameters lack direct symbolic alignment, this extraction procedure does not apply to the LB Transformer, which we therefore do not evaluate under applicability supervision.

% \paragraph{Direct extraction via symbolic alignment.}
%Under applicability supervision, traces contain their initial state but no subsequent states, so state reconstruction cannot be used.
%For the LS Transformer, however, the lifted \strips model can instead be read
%directly from its learned parameters. After binarization, the query parameters
%$q_{p,\alpha,\rho}$ specify whether the lifted atom given by
%$(p,\rho)$ is a precondition of $\alpha$, the key parameters
%$k_{p,\alpha,\rho}$ specify whether it is an effect, and the value parameters
%$v_{p,\alpha,\rho}$ specify the sign of that %effect, with $v=0$ for an add
%effect and $v=1$ for a delete effect. To associate the hidden predicate heads
%with the predicates of the ground-truth \strips model, we use the $\texttt{init-}p$ actions:
%$\texttt{init-}p$ has $p$ as its single add effect, and therefore identifies the hidden
% predicate corresponding to $p$. Hidden predicates for which no such alignment is available retain an anonymous name in the extracted model. In contrast, the parameters of the LB Transformer have no direct correspondence with \strips preconditions or effects. Although the LB Transformer could in principle be trained with applicability supervision, one cannot extract a symbolic model from it in this setting, so we do not consider this combination in our experiments.

\paragraph{Planning.}
Given a problem $P=\langle M,O,I,G\rangle$, we replace $M$ with the extracted model $M'$ and translate the initial state and goal using the learned predicate alignment. We use Mimir~\citep{mimir} with greedy best-first search and the FF heuristic~\citep{hoffmann2001ff}, limiting search to $10^7$ generated nodes. We validate each returned plan against ground-truth model $M$ and $P$.

% \paragraph{Planning.}
% Once a \strips model $M'$ has been extracted from the transformer, it can be used for planning with an off-the-shelf symbolic planner. Let $M$ denote the ground-truth \strips model, and $P$ a problem over $M$ which we want to solve. We can simply give $M'$ and $P$ as inputs to the planner. If $M'$ is equivalent to $M$, then the plan returned by the planer will be a valid solution for $P$. In our experiments, we use the Mimir planner~\citep{mimir} with greedy best-first search (GBFS) and the FastForward (FF) heuristic~\citep{hoffmann2001ff}.
% For each problem, search is given a budget of at most $10^7$ generated nodes.